\section{方法选择：为何不用 Atlas DCF}

Atlas 950 缺少客户验收数量、系统 ASP、公司 BOM 份额、收入确认、毛利率、营运资金和资本开支，无法建立可审计的产品 DCF。强行假设会把最重要的不确定性藏在一个精确目标价中。因此，完整模型以去除未桥接成长溢价的成熟业务 2026E PE 为主方法，以 FY2025 交易所年报归母净资产和标准化 ROE 推导的合理 P/B 为独立交叉验证；Atlas 专属 2026E 收入、净利和 EPS 均设为 0。所有 EPS 统一以“当前市值/收盘价”得到的现股本重算，券商表内 EPS 仅保留为送转股口径诊断。

每家公司目标价由基本面基准目标、当前市场情绪锚和原始卖方明示目标组成：
\[
P^{*}=w_f\times P_{\mathrm{base}}+w_m\times P_{\mathrm{market}}+w_b\times P_{\mathrm{broker}},\quad w_f+w_m+w_b=1.
\]
基本面目标由成熟业务 2026E PE 与合理 P/B 按 60\%/40\% 合成。合理 P/B 使用：
\[
P/B=\frac{ROE-g}{k_e-g}.
\]
Bear/Base/Bull 的股权成本为 11\%/10\%/9\%，永续增长为 2\%/3\%/4\%，ROE 上限为 23\%/25\%/27\%，P/B 下限为 0.6/0.7/0.8 倍。这些参数在 12 家公司间统一、事前固定，未用旧目标价反解。PEG 只保留为市场隐含预期诊断，目标价权重为 0；没有外部明示目标价时，$w_b=0$；市场锚通常只给 20\% 权重。Bear/Bull 同时改变 EPS/PE 和合理 P/B 参数，但三个情景均不计 Atlas 专属收入。

\begin{exhibitbox}[完整估值结果]
\scriptsize
\setlength{\tabcolsep}{1pt}
\begin{longtable}{L{1.4cm}L{2.0cm}R{1.4cm}R{1.4cm}R{1.4cm}R{1.4cm}R{1.4cm}R{1.5cm}R{1.5cm}}
\toprule
\textbf{代码} & \textbf{公司} & \textbf{现价} & \textbf{Bear} & \textbf{Base} & \textbf{Bull} & \textbf{综合目标} & \textbf{空间} & \textbf{2026E PE}\\
\midrule
\endhead
000034 & 神州数码 & 24.61 & 6.43 & 9.61 & 13.73 & 12.61 & -48.8\% & 27.7x\\
000988 & 华工科技 & 117.62 & 26.92 & 39.94 & 56.66 & 55.48 & -52.8\% & 49.8x\\
600183 & 生益科技 & 132.29 & 24.53 & 36.16 & 50.58 & 68.85 & -48.0\% & 57.7x\\
002916 & 深南电路 & 334.00 & 86.91 & 128.02 & 179.13 & 201.21 & -39.8\% & 41.6x\\
002463 & 沪电股份 & 127.80 & 30.55 & 45.54 & 64.06 & 61.99 & -51.5\% & 43.0x\\
300476 & 胜宏科技 & 241.50 & 85.75 & 129.90 & 183.65 & 152.22 & -37.0\% & 26.7x\\
002837 & 英维克 & 61.57 & 9.30 & 14.16 & 20.09 & 23.64 & -61.6\% & 68.6x\\
301018 & 申菱环境 & 86.96 & 11.34 & 17.64 & 25.84 & 31.50 & -63.8\% & 76.0x\\
002335 & 科华数据 & 29.75 & 8.09 & 11.54 & 15.98 & 15.18 & -49.0\% & 32.3x\\
002130 & 沃尔核材 & 15.24 & 10.26 & 15.50 & 22.46 & 15.45 & +1.4\% & 11.3x\\
002230 & 科大讯飞 & 41.12 & 6.87 & 9.73 & 12.92 & 16.01 & -61.1\% & 81.8x\\
002025 & 航天电器 & 64.76 & 9.11 & 14.11 & 20.31 & 42.06 & -35.1\% & 77.4x\\
\bottomrule
\end{longtable}
\sourcenote{价格日 2026-07-17；12 家交易所年报原件、38 份原始卖方报告、现股本复算与 PE/P\!B 双方法估值。目标只评价已披露业务，Atlas 专属收入信用为 0。综合目标仍按基本面、市场锚和明示卖方目标加权，因此个别公司可能高于 Base 情景。}
\end{exhibitbox}

\begin{exhibitbox}[完整最终估值矩阵：财务、方法与权重]
\scriptsize
\setlength{\tabcolsep}{1.5pt}
\begin{longtable}{L{1.0cm}R{1.25cm}R{1.25cm}R{1.05cm}R{1.0cm}R{1.15cm}R{1.15cm}L{1.25cm}R{1.15cm}R{1.1cm}L{1.7cm}}
\toprule
\textbf{代码} & \textbf{市值} & \textbf{26E收入} & \textbf{26E净利} & \textbf{EPS} & \textbf{PE目标} & \textbf{P/B目标} & \textbf{F/M/B} & \textbf{综合目标} & \textbf{空间} & \textbf{动作}\\
\midrule
\endhead
000034 & 250 & 1654.5 & 9.03 & 0.888 & 10.66 & 8.05 & 80/20/0 & 12.61 & -48.8\% & 观察\\
000988 & 1183 & 235.7 & 23.74 & 2.361 & 47.22 & 29.02 & 80/20/0 & 55.48 & -52.8\% & 观察\\
600183 & 3213 & 390.8 & 55.68 & 2.292 & 45.85 & 21.64 & 60/20/20 & 68.85 & -48.0\% & 避免追高\\
002916 & 2275 & 313.8 & 54.71 & 8.031 & 160.62 & 79.13 & 60/20/20 & 201.21 & -39.8\% & 观察\\
002463 & 2459 & 263.3 & 57.20 & 2.972 & 59.45 & 24.68 & 80/20/0 & 61.99 & -51.5\% & 观察\\
300476 & 2373 & 325.5 & 88.98 & 9.054 & 181.08 & 53.14 & 80/20/0 & 152.22 & -37.0\% & 选择性观察\\
002837 & 785 & 88.8 & 11.43 & 0.897 & 17.94 & 8.50 & 80/20/0 & 23.64 & -61.6\% & 等待毛利恢复\\
301018 & 325 & 67.5 & 4.28 & 1.143 & 20.58 & 13.21 & 80/20/0 & 31.50 & -63.8\% & 最高关系观察\\
002335 & 224 & 114.0 & 6.95 & 0.921 & 12.89 & 9.51 & 80/20/0 & 15.18 & -49.0\% & 等待估值回落\\
002130 & 213 & 121.5 & 18.80 & 1.343 & 16.12 & 14.59 & 80/20/0 & 15.45 & 1.4\% & 低纯度价值卫星\\
002230 & 988 & 329.0 & 12.07 & 0.502 & 12.56 & 5.48 & 80/20/0 & 16.01 & -61.1\% & 下游事件验证\\
002025 & 294 & 68.3 & 3.80 & 0.836 & 16.73 & 10.18 & 50/20/30 & 42.06 & -35.1\% & 关系未确认观察\\
\bottomrule
\end{longtable}
\sourcenote{单位：亿元、元。F/M/B 为基本面/市场/外部明示目标权重；PE 使用现股本重算 EPS，独立 P/B 锚使用交易所 FY2025 归母净资产、现股本、标准化 ROE、权益成本与永续增速，Atlas 专属收入、净利与 EPS 为 0。}
\end{exhibitbox}

\begin{exhibitbox}[完整最终估值矩阵：证据、催化与失效]
\scriptsize
\begin{longtable}{L{1.1cm}L{2.0cm}L{3.5cm}L{3.7cm}L{3.7cm}}
\toprule
\textbf{代码} & \textbf{证据等级} & \textbf{估值边界} & \textbf{下一验证/催化} & \textbf{失效/下调条件}\\
\midrule
\endhead
000034 & 广义业务可核验/Atlas 关系弱或未确认 & 已披露更广义业务；Atlas 信用为零 & 自有品牌毛利/现金转化；平台订单为额外升级项 & 分销利润或现金继续落后；路线图延期或无订单\\
000988 & 广义业务可核验/Atlas 关系弱或未确认 & 已披露更广义业务；Atlas 信用为零 & 800G+交付与光连接毛利；平台订单为额外升级项 & 交付或毛利低于预期；路线图延期或无订单\\
600183 & 广义业务可核验/Atlas 关系弱或未确认 & 已披露更广义业务；Atlas 信用为零 & 高端 CCL 结构/利用率；平台订单为额外升级项 & 产能爬坡或原料价差恶化；路线图延期或无订单\\
002916 & 广义业务可核验/Atlas 关系弱或未确认 & 已披露更广义业务；Atlas 信用为零 & AI PCB 良率/新线利用率；平台订单为额外升级项 & 高层板良率与毛利不达标；路线图延期或无订单\\
002463 & 广义业务可核验/Atlas 关系弱或未确认 & 已披露更广义业务；Atlas 信用为零 & AI PCB 出货/客户结构；平台订单为额外升级项 & 利用率或客户集中恶化；路线图延期或无订单\\
300476 & 广义业务可核验/Atlas 关系弱或未确认 & 已披露更广义业务；Atlas 信用为零 & 海外产能良率/AI 产品结构；平台订单为额外升级项 & 高增长和毛利不可持续；路线图延期或无订单\\
002837 & 广义业务可核验/Atlas 关系弱或未确认 & 已披露更广义业务；Atlas 信用为零 & Q2/H1 利润与毛利恢复；平台订单为额外升级项 & 利润恢复失败或订单弱；路线图延期或无订单\\
301018 & 关系较强/专属订单未证实 & 已披露更广义业务；Atlas 信用为零 & 数据服务订单转收入/现金；平台订单为额外升级项 & 订单不转收入或应收恶化；路线图延期或无订单\\
002335 & 广义业务可核验/Atlas 关系弱或未确认 & 已披露更广义业务；Atlas 信用为零 & 项目验收/扣非利润；平台订单为额外升级项 & 非经常收益主导或回款弱；路线图延期或无订单\\
002130 & 广义业务可核验/Atlas 关系弱或未确认 & 已披露更广义业务；Atlas 信用为零 & 高速线缆利用率/利润；平台订单为额外升级项 & 认证、产能或铜价传导失效；路线图延期或无订单\\
002230 & 关系较强/专属订单未证实 & 已披露更广义业务；Atlas 信用为零 & 10 月模型性能/付费转化；平台订单为额外升级项 & 发布延期或现金转化弱；路线图延期或无订单\\
002025 & 广义业务可核验/Atlas 关系弱或未确认 & 已披露更广义业务；Atlas 信用为零 & 平台认证/连接器收入；平台订单为额外升级项 & 关系仍未确认或军品周期下行；路线图延期或无订单\\
\bottomrule
\end{longtable}
\sourcenote{交易所年报原件、客户链审计、原始卖方盈利预测与风险矩阵。动作均为研究监测处置，不构成交易建议。}
\end{exhibitbox}


完整矩阵表明，方法分歧本身就是风险信息：成熟业务 PE 与合理 P/B 交叉目标差距较大时，说明当前盈利率、ROE 可持续性或资产效率需要更多验证；外部明示目标缺失时，券商目标权重严格为 0。动作表把每个目标价对应到下一季度催化与失效条件，避免只有数值而没有执行纪律。

\section{结果解读：没有同时满足“高关联+显著安全边际”的公司}

沃尔核材仅剩约 +1.4\% 的微弱目标空间，已不足以构成显著安全边际；其余 11 个完整模型均为负空间，胜宏科技也由此前的正空间改为约 -37.0\%。关系更强的申菱环境和科大讯飞，综合目标分别低于现价 63.8\% 和 61.1\%。换言之，最便宜的不是最“纯”，最“纯”的也不便宜。

这组结果不应解释成沃尔核材必然上涨或申菱环境必然下跌。它表达的是证据质量与当前价格的相对要求：低纯度公司需要靠原有业务兑现，关系强公司需要靠订单和利润快速追上估值。若 Atlas 订单得到确认，相关公司的 EPS 和倍数可以同时重估；在此之前，0 专属信用是更稳健的基线。

\section{相对估值、PEG 与 PSG 的适用边界}

所有完整模型均有正的 2026E EPS，因此成熟业务 PE 是主方法；合理 P/B 是独立二级锚，其官方净资产、BVPS、远期 ROE、股权成本和永续增长全部可复算。PEG 只适用于利润基数和口径相对稳定的公司；科大讯飞亏损季节性、拓维扣非亏损、航天电器盈利预测分歧等情形下，PEG 会制造失真。PSG 不用于目标价，因为分销、PCB、光器件、液冷和软件的毛利结构差异过大，同一收入增速不对应同一价值。

行业内部也不能用一个统一倍数：神州数码的分销和集成需要更低倍数并重视现金周转；光连接、AI PCB/CCL 的高端产品与客户认证可获得成长倍数，但需良率和产品结构验证；液冷需要订单、平台认证和毛利恢复；软件模型需要用户付费和经营杠杆。倍数是业务质量的结果，不是“华为链”标签的奖励。

\section{季节性与下一季度阈值}

模型没有把 2026Q1 机械年化。申菱环境与英维克 Q1 利润承压，而部分 PCB/CCL 公司 H1 预告强劲；科华数据 H1 归母和扣非增速差距较大；科大讯飞仍亏损。下一季度关键阈值分别是：申菱数据服务订单向收入和现金转化；英维克 Q2 毛利/利润恢复；华工光连接收入与毛利同步交付；科华项目验收和扣非利润质量；科大讯飞亏损收窄与付费商业化。

如果公司只披露“认证、送样、合作”，没有订单与收入，基本面目标不提高；如果披露订单但毛利率、应收和存货恶化，只提高收入不提高利润；如果收入、利润和现金流连续兑现，才同时上调 EPS 和估值倍数。

\section{市场隐含预期与压力测试}

以 2026E PE 作为市场隐含预期，样本从沃尔核材约 11 倍到科大讯飞约 82 倍；在 1.0 倍 PEG 的透明但非预测性约定下，这分别要求约 11\% 和 82\% 的可持续利润增长。若 EPS 下调 10\%、倍数压缩 20\%，公允价值约下降 28\%；若 EPS 上调 10\%、倍数扩张 20\%，约上升 32\%。高倍数公司的下行并不需要需求消失，只需增长久期或利润率低于预期。

\begin{riskbox}[估值审计结论]
估值模型复算：\textbf{通过}。价格、现股本、2026E/2027E 收入和净利、现股本重算 EPS、成熟业务 PE/合理 P\!B 双方法、Bear/Base/Bull、权重、目标价和上行空间均由结构化模型生成并独立复算；没有明示券商目标价的公司，Street 目标权重为 0；没有任何公司获得 Atlas 专属收入或利润信用。神州数码、英维克、申菱环境、科华数据等送转股口径差异较大的公司已不再直接使用券商表内 EPS。
\end{riskbox}
